\section{Validation, reproducibility, and exact scope}
\label{sec:validation}

\subsection{What was checked}
The finite correlation certificates and the rational schedule are replayable
with the Python standard library. The main checker exhausts every support
pair: 7,056 dot products for the four affine plans and 184,041 for the
thirteen cyclic plans. The six-group tail is checked as well. Integer power
comparisons certify the logarithm bounds; rational interval arithmetic
certifies the displayed constants and strict dominance of the depth-one
time branch. In particular the final intervals are
\[
 \begin{aligned}
 C_T&\in[30.413348726348439,30.413348726348440],\\
 C_S&\in[29.466030256594790,29.466030256594791].
 \end{aligned}
\]
An independently written checker reconstructs the plans and their supports
using bit sets, recomputes the integer certificates and rational packing
gates, and directly unrolls 93 levels of the paid recurrence as a numerical
cross-check. It does not import the primary checker. The finite unrolling
is supplemental: the proof handles the infinite geometric tail exactly.

The regrouping test uses 12 deterministic fixtures with at most 12 active
binary directed indices. It compares literal enumeration of global
assignments with separately implemented singleton caches, balanced binary
merges and the final three-block sum over
$(\mathbb Z/64\mathbb Z)[V]/(V^2)$. All 72 plan contractions agree.
The fixtures cover both orientations, spanning paths, disconnected
components, isolated scalar factors and zero factors. The nine fixtures
without a zero factor all have nonzero outputs. The checker also counts
every singleton entry call, verifies the count $\sum_v2^{\deg(v)}$,
and checks that the exact fifth-power selector minimizes the independently
reconstructed operation costs. This is a finite check of regrouping and
resource accounting, not an implementation of fast matrix multiplication.

For the local transform, 246 cases compare with independently formed
exact rational feature matrices, covering 10,160 entries; 27 larger cases
compare with the quadratic valuation-aware recurrence; and 96 odd-node
cases check the convolution specialization. Both signs, truncated mode
lists, missing ports, precision one, and the adjoint identity are covered.
The denominator guard bounds are checked in these cases. These finite
tests support the implementation; the product-tree and precision arguments
in Section~\ref{sec:transforms} establish the general identities.

All deterministic checks pass in normal Python and in \texttt{python -O},
with identical parsed output. Validation uses explicit exceptions rather
than depending on assertions that optimization could remove. The archived
results include both modes and a summary of the comparison.

\subsection{Matched local measurements}
The following scalar transpose measurements include preparation,
denominator inversions, increased working precision, and output recovery.
Every measured output was checked against the quadratic implementation.
Each number is the median of three runs at $p=4N$ and odd interpolation
valuation. The code uses CPython integers and Kronecker substitution;
these are small local measurements, not an end-to-end SAT benchmark.
\begin{center}\small
\begin{tabular}{r r r r r}
\toprule
$N$ & Quadratic (s) & One convolution (s) & Ratio & General tree (s)\\
\midrule
32 & 0.000407 & 0.000174 & $2.34\times$ & 0.002486\\
64 & 0.001923 & 0.000853 & $2.25\times$ & 0.011521\\
128 & 0.013941 & 0.007689 & $1.81\times$ & 0.091038\\
256 & 0.168149 & 0.057822 & $2.91\times$ & 0.656261\\
512 & 1.706872 & 0.418585 & $4.08\times$ & 6.034264\\
\bottomrule
\end{tabular}
\end{center}
The general product-tree prototype is slower at all these sizes. Its
asymptotic coefficient-operation bound remains useful, but these data
do not establish a practical crossover. The observed speedup belongs to
the odd-node specialization. The reference environment was Python 3.12.14
on Linux 6.18.55, x86\_64, glibc 2.39. Raw samples and environment metadata
are included so the measurements are not mistaken for portable constants.

\subsection{Proof status and limitations}
The newly proved parts are the correlation-certificate inequality, the
explicit affine and cyclic menus, their integration into a fully charged
four-stage recurrence, the precision-safe rational transforms and odd-node
convolution, and the restricted sparse-current obstruction. The result
uses the established interfaces of the inspected manuscripts: exact
Hermite extraction, normalized mode factorization, balanced packing,
entry semantics, SAT reduction, and the rational matrix multiplication
backend. Those dependencies are stated in Section~\ref{sec:interface};
this report is not an independent verification of every upstream claim.

The new contraction algorithm permits subexponential space. It does not
improve the strongest polynomial-space bound, prove a smaller logarithmic
order for unrestricted SAT, provide a finite-size crossover threshold,
or contain an end-to-end SAT implementation. Neither this new proof nor
its upstream dependencies were formalized or replayed in Lean here. The
source review of the supplied Anthropic revision does not change that
statement. The conditional graded-algebra proposition requires a new
construction meeting all of its hypotheses.

\subsection{Reproducing the archive}
From the extracted directory, run:
\begin{lstlisting}
python3 code/check_manifest.py
python3 code/run_checks.py
python3 code/build_report.py
\end{lstlisting}
The first command checks the delivered file hashes. The second replays
the finite certificates, independent audits, regrouping and kernel tests
in both Python modes. It requires only Python 3.10 or later. PDF rebuilding
requires pdfLaTeX and the packages named in \path{report.tex}; the supplied
figure is already present. To regenerate it, install Matplotlib and run
\path{code/make_figures.py}. Optional fresh timings use
\texttt{python3 code/fast\_cauchy\_transform.py --benchmark}; they need not
match the archived timings. The exact proof obligations, imported
dependencies and unproved targets are also recorded in
\path{PROOF_STATUS.json}.
